% --- IMPORTS ---
\documentclass[leqno]{article}
\usepackage{graphicx}
\usepackage{dsfont}
\usepackage{mathtools}
\usepackage{amssymb}
\usepackage{amsmath}
\usepackage{amsthm}
\usepackage{float}
\usepackage{ragged2e}
\usepackage{tensor}
\usepackage{fancyhdr}
\usepackage{empheq}
\usepackage[most]{tcolorbox}
\usepackage{enumitem}
\usepackage{dirtytalk}
\usepackage{marginnote}
\usepackage{microtype}
\usepackage[
  a4paper,
  left=0.8in,
  right=1in,
  top=1in,
  bottom=0.5in,
  headheight=14pt,
  headsep=0.4in
]{geometry}
\setlength{\parindent}{0cm}
% --- TikZ Packages ---
\usepackage{pgfplots}
\pgfplotsset{compat=1.18}
\usepgfplotslibrary{fillbetween, polar}
\usetikzlibrary{calc, positioning}
\usetikzlibrary{angles, quotes}
\usetikzlibrary{arrows.meta}
\usetikzlibrary{decorations.pathreplacing, decorations.markings}
\usetikzlibrary{patterns}
\usetikzlibrary{intersections}
% --- URL Packages ---
\usepackage[colorlinks=true,linkcolor=red,citecolor=magenta,urlcolor=black]{hyperref}%
\urlstyle{same}
% --- END OF IMPORTS ---

% --- CUSTOM COMMANDS ---
\RenewDocumentCommand{\marks}{o m}{%
  \IfValueTF{#1}%
    {\marginnote{\raggedright\textbf{[#2]}}[#1]}%
    {\marginnote{\raggedright\textbf{[#2]}}}%
}
\newcommand{\vspacerule}[2]{%
  \par\vspace*{#1}%
  \noindent\hrulefill\par
  \vspace*{#2}%
}
\definecolor{scarlet}{RGB}{205, 40, 75}
\newcommand{\questionitem}{%
  \item\phantomsection
  \addcontentsline{toc}{section}{Question \number\value{enumi}}%
}
\renewcommand{\contentsname}{Questions}
% --- END OF CUSTOM COMMANDS ---

% --- HEADER CONFIG ---
\pagestyle{fancy}
\fancyhead{}
\fancyfoot{}
\renewcommand{\headrulewidth}{0.9pt}
\lhead{Integration by Substitution (FM Level)}
\chead{}
\rhead{\href{https://danieljsmith.org}{danieljsmith.org}}
% --- END OF HEADER CONFIG ---

\begin{document}
\vspace*{20pt}
\tableofcontents
\newpage
\begin{enumerate}[
  label=\textbf{\arabic*.},
  leftmargin=*,
  topsep=0pt,
  partopsep=0pt,
  parsep=0pt
]

% ---- Question 1 ----
\questionitem
% [Source: _QBT__Integration_by_Substitution__FM_Level_, Q1]

Show that
\[
\int_0^{\pi/2} \frac{\cos^5 \theta}{2 + \sin \theta}\,\mathrm{d}\theta = 9\ln\!\left(\frac{3}{2}\right) - \frac{41}{12}
\]
\marks[-30pt]{8}

\newpage

% ---- Question 2 ----
\questionitem
% [Source: _QBT__Integration_by_Substitution__FM_Level_, Q2]

\begin{enumerate}[label=\textbf{(\alph*)}]
  \item Use the substitution $u = 1 + x^2$ to show that
  \[
  \int_0^2 \frac{x^9}{1+x^2}\,\mathrm{d}x = \int_p^q \frac{(u-1)^4}{2u}\,\mathrm{d}u
  \]
  where $p$ and $q$ are constants to be found. \marks{3} \\
  \item Hence show that
  \[
  \int_0^2 \frac{x^9}{1+x^2}\,\mathrm{d}x = A + B\ln 5
  \]
  where $A$ and $B$ are constants to be found. \marks{5} \\
\end{enumerate}

\newpage

% ---- Question 3 ----
\questionitem
% [Source: _QBT__Integration_by_Substitution__FM_Level_, Q3]

Show that 

\[
\int \frac{\sqrt{x^2 - 16}}{x^2}\,\mathrm{d}x = \operatorname{arcosh}\left(\frac{x}{4}\right) - \frac{\sqrt{x^2 - 16}}{x} + c
\]
\marks[-30pt]{6}

\newpage

% ---- Question 4 ----
\questionitem
% [Source: _QBT__Integration_by_Substitution__FM_Level_, Q4]

Find the exact value of 
\[
\int_0^{\frac{1}{2}\ln(2+\sqrt{3})} \frac{\operatorname{sech}^2 x}{\sqrt{4 - 9\tanh^2 x}}\,\mathrm{d}x
\]
\marks[-30pt]{8}

\newpage

% ---- Question 5 ----
\questionitem
% [Source: _QBT__Integration_by_Substitution__FM_Level_, Q5]

Using calculus, find the exact values of \\
\begin{enumerate}[label=(\roman*)]
  \item
  \[
  \int_1^2 \frac{1}{x^2 + 3x + 2}\,\mathrm{d}x
  \]
  \marks[-20pt]{3}
  \item
  \[
  \int_1^{\sqrt{3}} \frac{\sqrt{4 - x^2}}{x^2}\,\mathrm{d}x
  \]
  \marks[-30pt]{5}
\end{enumerate}

\newpage

% ---- Question 6 ----
\questionitem
% [Source: _QBT__Integration_by_Substitution__FM_Level_, Q6]

Find the exact value of \\
\[
\int_2^4 \frac{\sqrt{x^2-4}}{x^2} \, \mathrm{d}x
\]
Give your answer in simplest exact form. \marks{9}

\newpage

% ---- Question 7 ----
\questionitem
% [Source: _QBT__Integration_by_Substitution__FM_Level_, Q7]

Show that\\
\[
\int \frac{x^2}{\sqrt{x^2 + 16}}\,\mathrm{d}x = \frac{1}{2}x\sqrt{x^2 + 16} - 8\operatorname{arsinh}\left(\frac{x}{4}\right) + c
\]
\marks[-30pt]{6}

\newpage

% ---- Question 8 ----
\questionitem
% [Source: _QBT__Integration_by_Substitution__FM_Level_, Q8]

Show that
\[
\int_0^{\ln 2} \frac{e^{3x}+2e^{2x}}{\sqrt{e^x+4}}\,\mathrm{d}x = \frac{32\sqrt{6}-30\sqrt{5}}{5}
\]
\marks[-30pt]{7}

\newpage

% ---- Question 9 ----
\questionitem
% [Source: _QBT__Integration_by_Substitution__FM_Level_, Q9]

Using the substitution $u = \frac{2}{3}\tanh x$, evaluate
\[
\int_0^1 \frac{\operatorname{sech}^2 x}{\sqrt{9 - 4\tanh^2 x}}\,\mathrm{d}x
\]
\marks[-30pt]{8}

\newpage

% ---- Question 10 ----
\questionitem
% [Source: _QBT__Integration_by_Substitution__FM_Level_, Q10]

Show that
\[
\int_0^{\pi/2} \frac{\sin 3\theta}{2 + \cos \theta}\,\mathrm{d}\theta = 15\ln\left(\frac{3}{2}\right) - 6
\]
\marks[-20pt]{7}

\newpage

% ---- Question 11 ----
\questionitem
% [Source: _QBT__Integration_by_Substitution__FM_Level_, Q11]

\begin{enumerate}[label=\textbf{(\alph*)}]
  \item Use the substitution $u=\sqrt{2x+1}$ to show that
  \[
  \int_0^{15/2} \frac{x}{2+\sqrt{2x+1}}\,\mathrm{d}x = \int_p^q \frac{u(u^2-1)}{2(u+2)}\,\mathrm{d}u
  \]
  where $p$ and $q$ are constants to be found.
  \marks{3} \\
  \item Hence show that
  \[
  \int_0^{15/2} \frac{x}{2+\sqrt{2x+1}}\,\mathrm{d}x = A - B\ln 2
  \]
  where $A$ and $B$ are constants to be found.
  \marks{4}
\end{enumerate}

\newpage

% ---- Question 12 ----
\questionitem
% [Source: _QBT__Integration_by_Substitution__FM_Level_, Q12]

Use an appropriate substitution to find the exact value of
\[
\int_0^{\pi/6} \frac{\cos x\cos 2x}{(1 + \sin x)^{3/2}}\,\mathrm{d}x
\]
Give your answer in the form $a+b\sqrt{6}$, where $a$ and $b$ are rational numbers to be determined. \marks{7} \\

\newpage

% ---- Question 13 ----
\questionitem
% [Source: _QBT__Integration_by_Substitution__FM_Level_, Q13]

Using calculus, find the exact values of \\

\begin{enumerate}[label=(\roman*)]
  \item
  \[
  \int_{-2}^0 \frac{1}{x^2 + 4x + 8}\,\mathrm{d}x
  \]
  \marks[-30pt]{3}
  \item
  \[
  \int_0^{\sqrt{3}} \frac{x^2}{\sqrt{4 - x^2}}\,\mathrm{d}x
  \]
  \marks[-30pt]{5}
\end{enumerate}

\newpage

% ---- Question 14 ----
\questionitem
% [Source: _QBT__Integration_by_Substitution__FM_Level_, Q14]

Show that
\[
\int_1^3 \frac{(x-1)^2}{\sqrt{x+1}}\,\mathrm{d}x = \frac{16}{15}(7-4\sqrt{2})
\]
\marks[-20pt]{7}

\newpage

% ---- Question 15 ----
\questionitem
% [Source: _QBT__Integration_by_Substitution__FM_Level_, Q15]

Use an appropriate substitution to find the exact value of
\[
\int_0^{2\sqrt{2}} \frac{x^2}{\sqrt{16-x^2}}\,\mathrm{d}x
\]
Give your answer in the form $a+b\pi$, where $a$ and $b$ are rational numbers. \marks{9} \\
\end{enumerate}
\end{document}
